\chapter{Logistic Model}
\label{ch:08} 

\chaptermark{Chapter 8}

\abstract{This chapter presents the logistic growth model, a population dynamics model that incorporates density-dependent negative feedback on growth rates. After historically introducing population regulation concepts from Malthus and Verhulst's development of the first quantitative model, we derive the logistic differential equation and its closed-form solution. The model exhibits long-term convergence to a stable carrying capacity regardless of initial conditions. Through local stability analysis of steady states, analytical and graphical approaches confirm the model predicts sigmoidal growth trajectories for starts below half carrying capacity, and monotonic convergence above it. While density-dependence captures universal population constraints, the simplicity of the logistic equation demonstrates how fundamental processes can form the basis for mathematical descriptions reflecting real ecological dynamics over long timescales, making it one of the most widely applied population models.}

\vspace{1cm}

\section{Introduction}

While the exponential growth model provides a simple and mathematically tractable framework for studying populations over short time periods, it has significant limitations when extrapolated over long timescales. Most notably, the model predicts unconstrained population expansion leading to infinite size, which is not feasible in reality. All natural populations are bound by finite resources in their environment. The exponential model also ignores factors like resource competition and carrying capacity effects that act to slow growth as the population approaches the limits of its habitat. As a result, the model cannot realistically depict long-term population behavior or stability. For populations existing in a closed ecological system with density-dependent influences, another type of population model is needed that incorporates limitations on growth from resource limitations and crowding effects. This will allow an investigation of population regulation and equilibrium conditions. The development of such a model, which involves including density-dependent feedback in the growth term, is necessary to more accurately represent population dynamics over long periods aligned with biological timescales.

\section{A Brief Historical Review}

Thomas Malthus was one of the earliest scholars to study population dynamics. In his 1798 treatise \emph{An Essay on the Principle of Population}, Malthus hypothesized that human populations have an inherent capacity to grow exponentially if unchecked by outside forces. This is because populations can increase geometrically through reproduction. In contrast, he argued that the food resources required to support growing populations only increase in an arithmetic progression over time. As populations grow exponentially yet resources increase linearly, Malthus recognized that individuals would inevitably come into competition for limiting resources like food, space, and other necessities of life. This competition, he proposed, serves to slow population growth from its unchecked exponential rate down to a level sustainable given the available resources. Malthus' ideas introduced the concept that populations are regulated by external factors related to resource availability and by the scarcity and struggle for existence that arises due to overpopulation. His insights formed the basis for later formal mathematical models of population dynamics.

While Malthus recognized that populations cannot grow exponentially indefinitely, he did not provide a mathematical model to describe how growth becomes limited over time. It was the Belgian mathematician Pierre Verhulst who first developed an explicit equation incorporating this regulation mechanism. In 1838, Verhulst modified the simple exponential growth model based on Malthus' arguments about competition for resources. He proposed adding a negative feedback term that increases in proportion to the population size, representing the growing scarcity and resultant effect on per capita growth rate as the population approaches the environment's carrying capacity.

\section{The Logistic Model}

We begin by recalling the differential equation that governs exponential population growth:
\begin{equation*}
\frac{dN}{dt} = \alpha N,
\end{equation*}
where $\alpha$ is a constant growth rate parameter. Verhulst proposed modifying this model by substituting parameter $\alpha$ with a decreasing function of the population size $N$, denoted $\alpha(N)$. Specifically, he proposed a linear decay function of the form:
\begin{equation*}
\alpha(N) = r\left(1 - \frac{N}{K}\right),
\end{equation*}
where $r$ represents the intrinsic per capita growth rate in the absence of density-dependent effects. Meanwhile, $K$ denotes the environment's biological carrying capacity; the maximum sustainable population size given resource constraints. This negative feedback term grows proportionally as $N$ approaches $K$, introducing density-dependent regulation into the exponential growth formulation. Verhulst's ingenious substitution of a variable growth rate function was a landmark development that produced the first quantitative model embodying Malthus' concept of self-limitation through resource competition.

Through this modification, Verhulst produced the following nonlinear ordinary differential equation to govern population dynamics:
\begin{equation}
\frac{dN}{dt} = rN\left(1 - \frac{N}{K}\right).
\label{eq:08.01}
\end{equation}
Verhulst named the solution to this differential equation the ``logistic function". By extension, the overall model is commonly referred to as the ``logistic model" of population growth. However, it is also appropriately termed the ``Verhulst model"

We can solve the logistic differential equation (Eq. \ref{eq:08.01}) using variable transformations. First, introduce the change of variable $x = 1/N$. This transforms Eq. \eqref{eq:08.01} to:
\begin{equation*}
\frac{dx}{dt} = rx\left(\frac{1}{K} - x\right).
\end{equation*}
Next, define the variable $y = x - 1/K$, yielding:
\begin{equation*}
\frac{dy}{dt} = -ry.
\end{equation*}
This has the solution $y(t) = y_0e^{-rt}$, where $y_0$ is the initial value. Transforming back through the steps, we arrive at the solution for the original variable $N$:
\begin{equation}
N(t) = \left[\frac{1}{K} + \left(\frac{1}{N_0} - \frac{1}{K}\right)e^{-rt}\right]^{-1}.
\label{eq:08.02}
\end{equation}
Through judicious variable substitutions, we have derived an explicit closed-form solution to the logistic differential equation governing population dynamics. 

Eq. (\ref{eq:08.02}) is the well known logistic function. Consider the characteristics of this solution. The first term inside the brackets, $1/K$, represents the asymptotic carrying capacity. As time increases, the exponential term $e^{-rt}$ monotonically decays to zero. Therefore, in the limit of large time, the population $N$ will approach the value:
\begin{equation*}
\lim_{t\to\infty} N = K.
\end{equation*}
This shows that regardless of the initial population $N_0$, the solution will always converge to the environmental carrying capacity $K$. Physically, this makes sense; as time progresses, density-dependent effects will gradually dampen the growth rate until it balances perfectly with mortality to establish an equilibrium at the maximum sustainable size. The timescale to reach this steady state depends on the intrinsic growth rate $r$, but the ultimate attractor is always $K$. Therefore, the logistic model comprehensively captures the self-regulation of populations towards their local resource limits envisioned by Malthus.

\section{Local Stability of Steady States}

While we were able to derive an analytical solution to the logistic equation in the previous section, this is not always possible for nonlinear dynamical systems. Even when a closed-form solution cannot be obtained, important qualitative insights into dynamic behavior can still be gained by analyzing the system's steady states and their local stability properties. Though analytical solutions offer complete descriptions of trajectories over time, local stability analysis offers a general, mathematically rigorous technique to characterize global qualitative behaviors. In this section, we will use such an approach by finding steady states of the logistic model and applying linear stability theory to assess their stability.

To perform a local stability analysis, we first rewrite the logistic differential equation (Eq. \ref{eq:08.01}) in the general form:
\begin{equation*}
\frac{dN}{dt} = f(N),
\end{equation*}
where $f(N)=rN(1-N/K)$. The steady state solutions, also known as fixed points, correspond to values of $N$ where $f(N)=0$. One can readily show that the fixed points are $N_1^*=0$ and $N_2^*=K$.

We next calculate the derivative of $f(N)$ to determine the slope at each fixed point:
\begin{equation*}
f'(N) = r - \frac{2rN}{K}.
\end{equation*}
Evaluating $f'$ gives: $f'(N_1^*)=r>0$ and $f'(N_2^*)=-r<0$. This indicates that perturbations will grow near $N_1^*=0$, making it unstable. Meanwhile, near the carrying capacity $N_2^*=K$, perturbations will decay over time, signifying stability.

\begin{figure}
    \centering
    \includegraphics[width=3in]{"Figures/fig-08.01.pdf"}
    \caption{Graphical analysis of fixed points and stability in the logistic growth model. This figure shows the growth rate function $f(N)=rN(1-N/K)$. The two roots of this function, where it intersects the horizontal axis, are the fixed points (steady state solutions): $N_1^*=0$ and $N_2^*=K$. Local stability is determined by the slope of $f(N)$ at each fixed point. At $N_1^*=0$, the slope is positive indicating nearby trajectories will diverge over time. Therefore, $N_1^*$ denotes an unstable fixed point, represented by an open circle. In contrast, the slope of $f(N)$ is negative at the carrying capacity $N_2^*=K$. Thus, small perturbations will contract back towards the fixed point. This characterizes $N_2^*$ as a stable steady state solution, shown as a filled circle.}
    \label{fig:08.01}
\end{figure}

The previous analytic stability analysis can be visualized graphically. Fig. (\ref{fig:08.01}) depicts the growth rate function $f(N)=rN(1-N/K)$. This illustrates the two fixed points where $f(N)$ intersects the horizontal axis, representing steady state solutions.

The fixed points occur at population values of $N_1^*=0$ and $N_2^*=K$, coinciding with the steady state solutions found earlier. To assess local stability, we evaluate the slope of $f(N)$ at each fixed point. At $N_1^*=0$, the positive slope means nearby trajectories will diverge over time. Therefore, $N_1^*$ denotes an unstable node, shown as an open circle. In contrast, the negative slope at the carrying capacity $N_2^*=K$ implies small perturbations will decay back to this fixed point. This characterizes $N_2^*$ as locally stable, represented by a filled circle.

The graphical analysis in Fig. \ref{fig:08.01} provides additional qualitative insights into solution trajectories. Since the function $f(N)$ is increasing over the range $0<N<K/2$, populations initialized below half the carrying capacity ($N_0<K/2$) will experience accelerated growth as they approach this midpoint. However, as $f(N)$ is decreasing for $K/2<N<\infty$, these solutions will then decelerate as diminishing resources induce density-dependent regulation toward the asymptote. Thus, trajectories beginning below $K/2$ exhibit the characteristic S-shaped sigmoid curve.

On the other hand, solutions starting above half the carrying capacity ($N_0>K/2$) will converge directly to $K$ in a qualitatively different manner, as $f(N)$ is negative over this interval. Their approach will be steadily dampened regardless of whether their initial value is below or above the stable steady state $N_2^*=K$. Therefore, by examining the increasing and decreasing intervals of $f(N)$, we gain insight into qualitatively distinct population growth dynamics depending on initial population size relative to the system's threshold midpoint value.

\begin{figure}
    \centering
    \includegraphics[width=3in]{"Figures/fig-08.02.pdf"}
    \caption{Qualitative trajectories of the logistic solution for varying initial conditions. This figure graphically depicts the logistic function given by Eq. \eqref{eq:08.02} over time for different initial population sizes $N_0$. Solutions initialized below half the carrying capacity $K/2$ exhibit characteristic sigmoidal growth curves, accelerating as they approach the midpoint before slowing down due to density dependence. In contrast, populations starting above $K/2$ demonstrate qualitatively different monotonic convergence (without an inflection point) to the stable steady state at carrying capacity K, regardless of whether they start above or below the fixed point.}
    \label{fig:08.02}
\end{figure}

The qualitative insights gained from analyzing the growth rate function $f(N)$ can be validated by plotting the solutions of the logistic differential equation over time. Fig. \ref{fig:08.02} plots the logistic function given by Eq. \eqref{eq:08.02} for a range of initial population sizes $N_0$. This illustrates how initial conditions smaller than half the carrying capacity $K/2$ exhibit sigmoidal growth curves, consistent with the posited trajectory shapes based on the increasing and decreasing intervals of $f(N)$. Solutions initialized above $K/2$ also match expectations, showing direct monotonic convergence.

By comparing the predicted dynamics from analytic and graphical stability analysis (Fig. \ref{fig:08.01}) against analytic solutions in Fig. \ref{fig:08.02}, we find strong agreement between the local mathematical characterization near fixed points and the emergent global qualitative behaviors of the overall system across a spectrum of initial conditions. This corroborates that the linearization techniques employed to analyze steady states provide insights extending well beyond the local neighborhood of equilibria to inform our broader understanding of logistic population growth models.

\section{Discussion}

The logistic equation, despite its simple mathematical form, captures essential qualitative features of real-world population dynamics. The inclusion of nonlinear density-dependent feedback simulates the influence of finite resources and competition on growth. This confers carrying capacity behavior observed universally in populations, from microbial colonies to animal populations. Thus the logistic model has broad relevance in ecology, providing a starting point to understand outbreaks, booms, and stabilization in a variety of species. Its success stems from distilling dynamics to fundamental processes without species-specific complexity. As one of the earliest and most widely employed models in mathematical biology, it demonstrates how even minimal representations can reflect key characteristics of living populations, especially those of unicellular organisms like bacteria exhibiting exponential growth followed by limited expansion.

While analytic solutions offer complete trajectories when attainable, local stability analysis provides a generalizable technique for characterizing nonlinear dynamical systems. By linearizing near fixed points, it allows classifying long-term qualitative fates of solutions based on mathematical stability criteria. As demonstrated here for the logistic equation, this approach correctly predicts global dynamic attributes from investigation of a system's local equilibrium properties. Even without closed-form solutions, local analysis thus remains a powerful approach for divining qualitative insights into complexity arising from nonlinear feedbacks governing real natural and engineered systems. Where full descriptions remain elusive, studying stability properties of steady states provides an alternative lens into qualitative behaviors that may emerge from complexity near points of balance.

\section*{Exercises}

\begin{enumerate}
    \item Find the steady states of the following variation of the logistic model
    \begin{equation*}
        \frac{dN}{dt} = r N \left( 1 - \frac{N}{K}\right) - \delta.
    \end{equation*}
    Study the stability of steady states and discuss the implications in comparison with the behavior of the original model.

    \item Find the steady state and study their stability for the Gompertz model:
    \begin{equation*}
        \frac{dN}{dt} = r N \ln \frac{K}{N}.
    \end{equation*}
    Contrast the findings with those of the logistic model.
\end{enumerate}